% GENERATED FILE. Edit canonical lessons, then run npm run generate:books.
% canonical-source-hash: fnv1a64:5448ec936be8f621
% canonical-lessons: RU-C354-uluchshat, RU-C354-uvelichivatsya, RU-C354-umenshatsya, RU-C354-snizhat, RU-C354-povyshat

\chapter{Better or worse}
\label{ch:ru-better-or-worse}
\hlchaptermodality{354}

\begin{chapteropening}
\textbf{By the end of this chapter:} \emph{I can say that something improves or gets worse, grows or shrinks, and that someone lowers or raises it.}

Complete the last lesson of chapter 354: I can say that something improves or gets worse, grows or shrinks, and that someone lowers or raises it.
\end{chapteropening}

\section[\texorpdfstring{\ru{улучшать} (uluchshát)}{улучшать (uluchshát)}]{\ru{улучшать} (uluchshát) — to improve, to make better}

\label{lesson:RU-C354-uluchshat}



\begin{quote}
Before the new one: say the Russian for to suffer, then the Russian for to make worse.
\end{quote}

\subsection*{You'll want to know: \ru{улучшать}}
\textbf{\ru{улучшать}} (\emph{uluchshát}) — \textquotedblleft{}to improve, to make better\textquotedblright{}.

\begin{quote}
\ru{Чтение} \ru{улучшает} \ru{память}.
\end{quote}

\begin{quote}
\emph{Chténiye uluchsháyet pámyat.}
\end{quote}

\begin{quote}
Reading improves your memory.
\end{quote}

Imperfective; perfective \emph{\ru{улучшить}}. \emph{\ru{Улучшаться}} — to get better.

\begin{cousinweb}
\textbf{\ru{улучшать}} is built on \emph{\ru{лучший}} \textquotedblleft{}better\textquotedblright{}, which is usually traced to Old Russian \emph{\ru{лучии}} \textquotedblleft{}better, more fitting\textquotedblright{}.
\end{cousinweb}

\begin{grammarlens}[title={What you've built}]
One more word for telling, explaining and describing at B1.
\end{grammarlens}

\subsection*{Your turn}
\begin{itemize}
\raggedright
  \item \emph{Say it:} \emph{uluchshát}
  \item \emph{Say it:} \emph{uluchshát}, once more
  \item \emph{Say it:} a sentence of your own with \emph{uluchshát}
  \item \emph{Recall:} say the Russian for to suffer, then the Russian for to make worse, then say \emph{uluchshát} again
\end{itemize}

\begin{tcolorbox}[breakable,colback=teal!4,colframe=teal!35!black,title={Before you move on}]
What does \ru{улучшать} mean? (\textquotedblleft{}To improve, to make better\textquotedblright{}.) Say it once more.
\end{tcolorbox}
\section[\texorpdfstring{\ru{увеличиваться} (uvelíchivatsya)}{увеличиваться (uvelíchivatsya)}]{\ru{увеличиваться} (uvelíchivatsya) — to increase, to grow}

\label{lesson:RU-C354-uvelichivatsya}



\begin{quote}
Before the new one: say the Russian for to make worse, then the Russian for to improve.
\end{quote}

\subsection*{You'll want to know: \ru{увеличиваться}}
\textbf{\ru{увеличиваться}} (\emph{uvelíchivatsya}) — \textquotedblleft{}to increase, to grow\textquotedblright{}.

\begin{quote}
\ru{Цены} \ru{увеличиваются} \ru{каждый} \ru{год}.
\end{quote}

\begin{quote}
\emph{Tsény uvelíchivayutsya kázhdyy god.}
\end{quote}

\begin{quote}
Prices increase every year.
\end{quote}

Imperfective; the perfective is \emph{\ru{увеличиться}}. The transitive verb is \emph{\ru{увеличивать}}, to increase something.

\begin{cousinweb}
\textbf{\ru{увеличиваться}} is built on \emph{\ru{великий}} \textquotedblleft{}great\textquotedblright{}, inherited from Proto-Slavic \emph{*velikŭ}, from \emph{*velĭjĭ} \textquotedblleft{}great\textquotedblright{} (EtymDB): to become greater.
\end{cousinweb}

\begin{grammarlens}[title={What you've built}]
One more word for telling, explaining and describing at B1.
\end{grammarlens}

\subsection*{Your turn}
\begin{itemize}
\raggedright
  \item \emph{Say it:} \emph{uvelíchivatsya}
  \item \emph{Say it:} \emph{uvelíchivatsya}, once more
  \item \emph{Say it:} a sentence of your own with \emph{uvelíchivatsya}
  \item \emph{Recall:} say the Russian for to make worse, then the Russian for to improve, then say \emph{uvelíchivatsya} again
\end{itemize}

\begin{tcolorbox}[breakable,colback=teal!4,colframe=teal!35!black,title={Before you move on}]
What does \ru{увеличиваться} mean? (\textquotedblleft{}To increase, to grow\textquotedblright{}.) Say it once more.
\end{tcolorbox}
\section[\texorpdfstring{\ru{уменьшаться} (umenshátsya)}{уменьшаться (umenshátsya)}]{\ru{уменьшаться} (umenshátsya) — to decrease, to get smaller}

\label{lesson:RU-C354-umenshatsya}



\begin{quote}
Before the new one: say the Russian for to improve, then the Russian for to increase.
\end{quote}

\subsection*{You'll want to know: \ru{уменьшаться}}
\textbf{\ru{уменьшаться}} (\emph{umenshátsya}) — \textquotedblleft{}to decrease, to get smaller\textquotedblright{}.

\begin{quote}
\ru{Боль} \ru{постепенно} \ru{уменьшается}.
\end{quote}

\begin{quote}
\emph{Bol postepénno umensháyetsya.}
\end{quote}

\begin{quote}
The pain is gradually decreasing.
\end{quote}

Imperfective; perfective \emph{\ru{уменьшиться}}. Built on \emph{\ru{меньше}}, less.

\begin{cousinweb}
\textbf{\ru{уменьшаться}} is built from \emph{\ru{у}-} and \emph{\ru{меньше}} \textquotedblleft{}less\textquotedblright{}, an old comparative Russian inherited from Proto-Slavic; it is not built on \emph{\ru{малый}}, which has a different root.
\end{cousinweb}

\begin{grammarlens}[title={What you've built}]
One more word for telling, explaining and describing at B1.
\end{grammarlens}

\subsection*{Your turn}
\begin{itemize}
\raggedright
  \item \emph{Say it:} \emph{umenshátsya}
  \item \emph{Say it:} \emph{umenshátsya}, once more
  \item \emph{Say it:} a sentence of your own with \emph{umenshátsya}
  \item \emph{Recall:} say the Russian for to improve, then the Russian for to increase, then say \emph{umenshátsya} again
\end{itemize}

\begin{tcolorbox}[breakable,colback=teal!4,colframe=teal!35!black,title={Before you move on}]
What does \ru{уменьшаться} mean? (\textquotedblleft{}To decrease, to get smaller\textquotedblright{}.) Say it once more.
\end{tcolorbox}
\section[\texorpdfstring{\ru{снижать} (snizhát)}{снижать (snizhát)}]{\ru{снижать} (snizhát) — to lower, to reduce}

\label{lesson:RU-C354-snizhat}



\begin{quote}
Before the new one: say the Russian for to increase, then the Russian for to decrease.
\end{quote}

\subsection*{You'll want to know: \ru{снижать}}
\textbf{\ru{снижать}} (\emph{snizhát}) — \textquotedblleft{}to lower, to reduce\textquotedblright{}.

\begin{quote}
\ru{Магазин} \ru{снижает} \ru{цены} \ru{перед} \ru{праздниками}.
\end{quote}

\begin{quote}
\emph{Magazín snizháyet tsény péred prázdnikami.}
\end{quote}

\begin{quote}
The shop lowers its prices before the holidays.
\end{quote}

Imperfective; perfective \emph{\ru{снизить}} (\emph{\ru{снижу}, \ru{снизишь}}). \emph{\ru{Снижаться}} — to go down.

\begin{cousinweb}
\textbf{\ru{снижать}} is built from \emph{\ru{с}-} and \emph{\ru{низ}} \textquotedblleft{}the bottom\textquotedblright{}, from Proto-Slavic \emph{*nizŭ}, the root of \emph{\ru{низкий}} \textquotedblleft{}low\textquotedblright{} (EtymDB).
\end{cousinweb}

\begin{grammarlens}[title={What you've built}]
One more word for telling, explaining and describing at B1.
\end{grammarlens}

\subsection*{Your turn}
\begin{itemize}
\raggedright
  \item \emph{Say it:} \emph{snizhát}
  \item \emph{Say it:} \emph{snizhát}, once more
  \item \emph{Say it:} a sentence of your own with \emph{snizhát}
  \item \emph{Recall:} say the Russian for to increase, then the Russian for to decrease, then say \emph{snizhát} again
\end{itemize}

\begin{tcolorbox}[breakable,colback=teal!4,colframe=teal!35!black,title={Before you move on}]
What does \ru{снижать} mean? (\textquotedblleft{}To lower, to reduce\textquotedblright{}.) Say it once more.
\end{tcolorbox}
\section[\texorpdfstring{\ru{повышать} (povyshát)}{повышать (povyshát)}]{\ru{повышать} (povyshát) — to raise, to increase}

\label{lesson:RU-C354-povyshat}



\begin{quote}
Before the new one: say the Russian for to decrease, then the Russian for to lower.
\end{quote}

\subsection*{You'll want to know: \ru{повышать}}
\textbf{\ru{повышать}} (\emph{povyshát}) — \textquotedblleft{}to raise, to increase\textquotedblright{}.

\begin{quote}
\ru{Кофе} \ru{повышает} \ru{давление}.
\end{quote}

\begin{quote}
\emph{Kófe povysháyet davléniye.}
\end{quote}

\begin{quote}
Coffee raises blood pressure.
\end{quote}

Imperfective; the perfective is \emph{\ru{повысить}} (\emph{\ru{повышу}, \ru{повысишь}}). \emph{\ru{Повышать} \ru{голос}} — to raise one’s voice.

\begin{cousinweb}
\textbf{\ru{повышать}} is built from \emph{\ru{по}-} and \emph{\ru{выше}} \textquotedblleft{}higher\textquotedblright{}, the comparative of \emph{\ru{высокий}} \textquotedblleft{}high\textquotedblright{}, inherited from Proto-Slavic \emph{*vysokŭ} \textquotedblleft{}high, tall\textquotedblright{} (EtymDB).
\end{cousinweb}

\begin{grammarlens}[title={What you've built}]
One more word for telling, explaining and describing at B1.
\end{grammarlens}

\subsection*{Your turn}
\begin{itemize}
\raggedright
  \item \emph{Say it:} \emph{povyshát}
  \item \emph{Say it:} \emph{povyshát}, once more
  \item \emph{Say it:} a sentence of your own with \emph{povyshát}
  \item \emph{Recall:} say the Russian for to decrease, then the Russian for to lower, then say \emph{povyshát} again
\end{itemize}

\begin{tcolorbox}[breakable,colback=teal!4,colframe=teal!35!black,title={Before you move on}]
What does \ru{повышать} mean? (\textquotedblleft{}To raise, to increase\textquotedblright{}.) Say it once more.
\end{tcolorbox}
